\section{Selecting the number of Kira workers}
\label{app:thread-policy}
The worker-count study retains the seed strategy from Section~\ref{sec:basis-completion} and varies only the number of Kira workers used for a requested batch. The families and target patterns are those in Appendix~\ref{app:type-specialists}. Parameter profile 91 supplies the
development test, profiles 92 and 93 supply training, profile 94 supplies
validation, and profiles 95 and 96 supply the reserved comparison.
The dimension remains symbolic. Target patterns overlap across these
sets; the rational kinematic values differ.

\subsection{Choices and exactness}
The preceding learned method chooses seed bounds and can reduce through
a temporary basis. We retain that method and let a new policy choose
one, two, or four Kira workers for each requested batch. The choice
applies to each Kira attempt, including a temporary-basis conversion.
Kira supports parallel symbolic reduction~\cite{KiraOriginal}. Changing the worker count changes how operations are scheduled, while each operation continues to use exact rational arithmetic. Every output must reach the
prescribed master basis and agree with an independent expanded-seed
reduction. If a new choice fails, its time and the repeat with the
preceding one-worker method both enter the complete time.

Let $A\mathbf I=0$ denote the IBP system. An exact elimination step
left-multiplies it by an invertible elementary matrix $E$, so
\begin{equation}
 A\mathbf I=0\quad\Longleftrightarrow\quad
 EA\mathbf I=0.
 \label{eq:thread-row-equivalence}
\end{equation}
Multiplication by $E$ proves the forward implication, and multiplication by $E^{-1}$ proves the reverse. Repeating valid steps therefore preserves the system, provided parallel execution respects their dependencies. Kira manages those dependencies through its existing workers, while the predictor sets only their number.
This argument preserves relations but does not prove that a finite seed
set is sufficient. The basis and coefficient checks therefore remain
part of every measured reduction.

\subsection{Cost labels and selection}
Write $x$ for the nine target features defined in
Appendix~\ref{app:basis-completion}, and $q\in\{1,2,4\}$ for the worker
count. For training batch $b$, let $t_{b,q}$ be the median complete time of
two measurements. The training label is
\begin{equation}
 y_{b,q}=\log(t_{b,q}/t_{b,1}).
 \label{eq:thread-label}
\end{equation}
Because the logarithm is strictly increasing and $t_{b,1}>0$ is common
to all choices for that batch, minimizing this label also minimizes
$t_{b,q}$. We fit the same ten predictor types per graph as in
Appendix~\ref{app:basis-completion}. They include constants, trees,
bagged trees, regularized linear fits, and nearest-neighbour fits.

The training-selected fixed control minimizes $\sum_b t_{b,q}$.
Validation uses three measurements per choice and batch. It first selects
the faster of that fixed control and the preceding method, using the sum
of batch medians. It then compares each predictor with decision margins
$0$, $0.05$, and $0.15$ in log relative cost. A predictor is enabled only
if its summed median cost is at least two percent lower than the control
and its summed cost is lower in at least two of the three validation
repetitions. Otherwise the control stays active. This objective measures the time for the workload by summing batch costs. It therefore gives a different weighting from the geometric mean that treats graphs equally. The selected controls and models remain fixed throughout the reserved tests.
Unknown graph structures or features outside the training range use the
preceding method.

\subsection{Complete-time comparisons}
The reserved comparison has 36 batches and five repetitions. It measures
native choices with one worker, the preceding learned method with one
worker, that same method with four workers throughout, a training-selected
fixed worker count per graph, and the final selected method. All methods use the same integrated executable and its common earlier
optimizations. The native control uses Kira's native seed and reduction
choices in that executable. All methods run sequentially in the same
four-core allocation. Execution order rotates
between repetitions. Thus the always-four-worker control measures the
cost of using the extra cores without target-dependent selection.

Full time includes model decisions, input preparation, Kira startup,
equation generation, exact algebra, any basis conversion, coefficient
checks, fallback, and cleanup. The caller and models are already loaded,
while every attempt starts a new Kira process. Offline reference
construction, training, and record writing are excluded. CPU time sums
user and system time of the caller and its completed child processes.
This measures the CPU time actually consumed. Every method has four reserved cores, including the methods that use only one worker. The comparison reports both the sum of batch medians and
the graph-weighted geometric ratios defined in the earlier appendices.
Additional applications receive the targets and basis without reference
coefficients and are checked after reduction.
